\section{模型评价与推广}

\subsection{模型优点}
\begin{enumerate}[label=\arabic*.,leftmargin=2.4em,itemsep=0.35em]
  \item \textbf{模型体系统一且递进。}本文以模块旋转、坐标和非重叠关系为共同几何基础，从自由轮廓依次扩展至固定轮廓与网络线长、紧轮廓搜索，再将矩形表示推广至非矩形的真实占用区域集合。四个问题可复用共同的变量、约束与检验逻辑，同时保留各问的特有目标。
  \item \textbf{优化层级与题目目标一致。}问题一对外接面积和轮廓长短边比采用词典序判优，问题三则以“先压缩轮廓、再优化 HPWL”构造词典序目标。这一处理不需要人为设定面积与线长的加权系数，避免了额外权重对目标优先级的改变。
  \item \textbf{网络信息与离散几何约束分层处理。}问题二先由二次线长松弛提取网络驱动的目标位置，再进行几何合法化和局部重插入。这一流程将网络引导与几何可行性分工，缓解了位置、旋转、非重叠与 HPWL 在高维离散空间中直接联合搜索的困难。
  \item \textbf{结果具有可复现的证据链。}问题一至问题三均输出模块坐标、旋转状态与布局图，并独立复核完整性、边界、非重叠与目标函数；问题四进一步以理论面积下界和具体可行构造闭合最优性证明，使数值结果与结论强度保持匹配。
  \item \textbf{形状级几何表示能够利用非矩形凹槽。}问题四不以外接矩形代替 L 型、T 型模块，而依据真实占用集合判断碰撞，因而允许其他模块进入未被实体占用的凹槽。最终零死区布局说明，该表示能够避免外接矩形近似引起的虚假空间占用。
\end{enumerate}

\subsection{模型局限}
\begin{enumerate}[label=\arabic*.,leftmargin=2.4em,itemsep=0.35em]
  \item \textbf{大规模实例尚无法证明全局最优。}问题一的 Skyline 与 MaxRects 构造、问题二的目标引导合法化与局部重插入，以及问题三以 MaxRects 进行的当前轮廓可行构造，均只探索了组合空间的一部分。因此，相应面积、HPWL 和成功构造边长应解释为当前策略与搜索预算下的经验证可行上界，而非严格全局最优值。
  \item \textbf{HPWL 仅是实际互连代价的近似。}HPWL 能以较低代价衡量网络的几何展布范围，是 VLSI 物理设计中常用的线长估计指标~\cite{kahng2011}；但真实布线还受拥塞、障碍、布线层与时序等因素影响，因而较小 HPWL 不必然对应详细布线后的最小实际线长或最小时延。
  \item \textbf{引脚位置采用模块中心近似。}问题二与问题三将硬模块引脚统一置于几何中心。真实宏模块的引脚可能偏离中心且随旋转变化；当模块尺寸较大或引脚集中在某一侧时，中心近似可能改变 HPWL 对位置与旋转状态的评价。
  \item \textbf{问题四的完整枚举难以直接扩展。}当前实例仅含四个模块，旋转与整数位置可以被完整枚举。随着模块数量、轮廓尺度和形状复杂度增长，组合数量将迅速扩张，需要引入具备形状级碰撞检测的启发式或元启发式搜索。
  \item \textbf{部分结果依赖题意边界口径。}Terminal 是否必须位于芯片轮廓内，会改变问题三报告的死区率。本文的敏感性分析表明，该差异不改变主要结论；但在工程应用中，仍需根据芯片 I/O、固定宏单元与封装约束明确边界定义。
\end{enumerate}

\subsection{模型推广}
\begin{enumerate}[label=\arabic*.,leftmargin=2.4em,itemsep=0.35em]
  \item \textbf{引入结构化布图编码以扩大全局搜索范围。}对大规模矩形模块，可进一步采用 Sequence-Pair~\cite{murata1996} 或 B$^*$-Tree~\cite{chang2000} 等非切片布图表示，并结合模拟退火、多起点搜索或混合整数规划。相对单纯的顺序构造，结构化编码能显式表示模块间的相对位置，为更广泛的邻域搜索提供基础。
  \item \textbf{从单一 HPWL 扩展至多物理指标协同优化。}实际 VLSI 物理设计还需同时考虑拥塞、关键路径时延、功耗和热分布等指标~\cite{kahng2011}，可建立
  \begin{equation}
    \min\left(A,\mathcal{L}_{\mathrm{HPWL}},C_{\mathrm{congestion}},T_{\mathrm{delay}},P\right)
  \end{equation}
  的多目标模型。当各指标不存在严格优先关系时，可用 Pareto 前沿代替本文的词典序判优，为设计者提供面积、线长与性能之间的多种折中方案。
  \item \textbf{引入真实引脚位置与网络权重。}若能获得模块内部引脚的局部坐标，可将其随模块平移和旋转同步变换；同时依据时序关键程度设置网络权重，将普通 HPWL 推广为加权 HPWL，以提高模型与后续详细布局及布线代价的一致性。
  \item \textbf{推广至一般直角多边形模块。}对更复杂的 L 型、T 型与一般 rectilinear 模块，可继续以有限矩形并或多边形表示真实占用区域，并将形状级相交判定嵌入 B$^*$-Tree 等全局搜索结构中，使问题四的几何思想不再局限于整数单位格。
  \item \textbf{面向超大规模网表进行分层与并行优化。}当模块数量进一步增长时，可先按网络连接强度将模块聚类为超节点，完成粗粒度全局区域划分，再在各区域内进行精细布局；不同候选轮廓、模块顺序和局部搜索过程也可并行执行，以提高现有分层优化框架对更大规模实例的处理能力。
\end{enumerate}
